
\documentclass{article}
\usepackage{colortbl}
\usepackage{makecell}
\usepackage{multirow}
\usepackage{supertabular}

\begin{document}

\newcounter{utterance}

\centering \large Interaction Transcript for game `matchit\_ascii', experiment `similar\_grid\_2', episode 3 with Qwen3.6{-}35B{-}A3B{-}FP8{-}without{-}reasoning.
\vspace{24pt}

{ \footnotesize  \setcounter{utterance}{1}
\setlength{\tabcolsep}{0pt}
\begin{supertabular}{c@{$\;$}|p{.15\linewidth}@{}p{.15\linewidth}p{.15\linewidth}p{.15\linewidth}p{.15\linewidth}p{.15\linewidth}}
    \# & $\;$A & \multicolumn{4}{c}{Game Master} & $\;\:$B\\
    \hline

    \theutterance \stepcounter{utterance}  
    & & \multicolumn{4}{p{0.6\linewidth}}{
        \cellcolor[rgb]{0.9,0.9,0.9}{
            \makecell[{{p{\linewidth}}}]{
                \texttt{\tiny{[P1$\langle$GM]}}
                \texttt{Stai partecipando a un gioco di indovinelli collaborativo.} \\
\texttt{L'obiettivo è scoprire se questa griglia e un'altra immagine che vedo solo io sono uguali.} \\
\texttt{Per favore descrivi prima la tua griglia. Poi, io fornirò la mia descrizione e potremo farci domande reciproche sulle griglie per capire se sono uguali. Ora inizia la tua breve descrizione della griglia con "DESCRIZIONE:" seguito dalla descrizione. Non aggiungere altro. Questa è la tua griglia:} \\
\texttt{X X X X X} \\
\texttt{▢ ▢ ▢ ▢ X} \\
\texttt{▢ ▢ X X X} \\
\texttt{▢ ▢ ▢ ▢ X} \\
\texttt{▢ ▢ ▢ ▢ X} \\
            }
        }
    }
    & \\ \\

    \theutterance \stepcounter{utterance}  
    & \multicolumn{4}{p{0.6\linewidth}}{
        \cellcolor[rgb]{0.9,0.9,0.9}{
            \makecell[{{p{\linewidth}}}]{
                \texttt{\tiny{[P1$\rangle$GM]}}
                \texttt{DESCRIZIONE: La griglia è una matrice 5x5 composta da due simboli: "X" e "▢". La prima riga è interamente composta da "X". La seconda riga ha quattro "▢" seguiti da una "X". La terza riga ha due "▢" seguiti da tre "X". La quarta e la quinta riga sono identiche, composte da quattro "▢" seguiti da una "X".} \\
            }
        }
    }
    & & \\ \\

    \theutterance \stepcounter{utterance}  
    & & & \multicolumn{2}{p{0.3\linewidth}}{
        \cellcolor[rgb]{0.9,0.9,0.9}{
            \makecell[{{p{\linewidth}}}]{
                \texttt{\tiny{[GM$|$GM]}}
                \texttt{continue} \\
            }
        }
    }
    & & \\ \\

    \theutterance \stepcounter{utterance}  
    & & \multicolumn{4}{p{0.6\linewidth}}{
        \cellcolor[rgb]{0.9,0.9,0.9}{
            \makecell[{{p{\linewidth}}}]{
                \texttt{\tiny{[GM$\rangle$P2]}}
                \texttt{Stai partecipando a un gioco di indovinelli collaborativo.} \\
\texttt{L'obiettivo è scoprire se questa griglia e un'altra immagine che vedo solo io sono uguali.} \\
\texttt{Per favore descrivi prima la tua griglia. Poi, io fornirò la mia descrizione e potremo farci domande reciproche sulle griglie per capire se sono uguali. Ora inizia la tua breve descrizione della griglia con "DESCRIZIONE:" seguito dalla descrizione. Non aggiungere altro. Questa è la tua griglia:} \\
\texttt{X X X X X} \\
\texttt{▢ ▢ ▢ ▢ X} \\
\texttt{▢ X X X X} \\
\texttt{▢ ▢ ▢ ▢ X} \\
\texttt{▢ ▢ X ▢ X} \\
            }
        }
    }
    & \\ \\

    \theutterance \stepcounter{utterance}  
    & & & \multicolumn{4}{p{0.6\linewidth}}{
        \cellcolor[rgb]{0.9,0.9,0.9}{
            \makecell[{{p{\linewidth}}}]{
                \texttt{\tiny{[GM$\langle$P2]}}
                \texttt{DESCRIZIONE: La griglia è una matrice 5x5 composta da due simboli: "X" e "▢".} \\
\texttt{Riga 1: X X X X X} \\
\texttt{Riga 2: ▢ ▢ ▢ ▢ X} \\
\texttt{Riga 3: ▢ X X X X} \\
\texttt{Riga 4: ▢ ▢ ▢ ▢ X} \\
\texttt{Riga 5: ▢ ▢ X ▢ X} \\
            }
        }
    }
     \\ \\

    \theutterance \stepcounter{utterance}  
    & & & \multicolumn{2}{p{0.3\linewidth}}{
        \cellcolor[rgb]{0.9,0.9,0.9}{
            \makecell[{{p{\linewidth}}}]{
                \texttt{\tiny{[GM$|$GM]}}
                \texttt{continue} \\
            }
        }
    }
    & & \\ \\

    \theutterance \stepcounter{utterance}  
    & & \multicolumn{4}{p{0.6\linewidth}}{
        \cellcolor[rgb]{0.9,0.9,0.9}{
            \makecell[{{p{\linewidth}}}]{
                \texttt{\tiny{[GM$\rangle$P2]}}
                \texttt{Questa è la mia DESCRIZIONEDESCRIZIONE: La griglia è una matrice 5x5 composta da due simboli: "X" e "▢". La prima riga è interamente composta da "X". La seconda riga ha quattro "▢" seguiti da una "X". La terza riga ha due "▢" seguiti da tre "X". La quarta e la quinta riga sono identiche, composte da quattro "▢" seguiti da una "X".} \\
\texttt{Ora fai una domanda per scoprire nuovi aspetti della mia griglia che potrebbero essere diversi dalla tua griglia. Inizia con "DOMANDA:" e non aggiungere altro.} \\
            }
        }
    }
    & \\ \\

    \theutterance \stepcounter{utterance}  
    & & & \multicolumn{4}{p{0.6\linewidth}}{
        \cellcolor[rgb]{0.9,0.9,0.9}{
            \makecell[{{p{\linewidth}}}]{
                \texttt{\tiny{[GM$\langle$P2]}}
                \texttt{DOMANDA: Nella tua descrizione, affermi che la quarta e la quinta riga sono identiche e composte da quattro "▢" seguiti da una "X". Nella mia griglia, la quarta riga è effettivamente "▢ ▢ ▢ ▢ X", ma la quinta riga è "▢ ▢ X ▢ X". Puoi confermare se nella tua griglia la quinta riga è davvero uguale alla quarta, o se c'è una differenza nella posizione della "X" nell'ultima riga?} \\
            }
        }
    }
     \\ \\

    \theutterance \stepcounter{utterance}  
    & & & \multicolumn{2}{p{0.3\linewidth}}{
        \cellcolor[rgb]{0.9,0.9,0.9}{
            \makecell[{{p{\linewidth}}}]{
                \texttt{\tiny{[GM$|$GM]}}
                \texttt{continue} \\
            }
        }
    }
    & & \\ \\

    \theutterance \stepcounter{utterance}  
    & & \multicolumn{4}{p{0.6\linewidth}}{
        \cellcolor[rgb]{0.9,0.9,0.9}{
            \makecell[{{p{\linewidth}}}]{
                \texttt{\tiny{[P1$\langle$GM]}}
                \texttt{Questa è la mia DESCRIZIONEDESCRIZIONE: La griglia è una matrice 5x5 composta da due simboli: "X" e "▢".} \\
\texttt{Riga 1: X X X X X} \\
\texttt{Riga 2: ▢ ▢ ▢ ▢ X} \\
\texttt{Riga 3: ▢ X X X X} \\
\texttt{Riga 4: ▢ ▢ ▢ ▢ X} \\
\texttt{Riga 5: ▢ ▢ X ▢ X} \\
\texttt{DOMANDA: Nella tua descrizione, affermi che la quarta e la quinta riga sono identiche e composte da quattro "▢" seguiti da una "X". Nella mia griglia, la quarta riga è effettivamente "▢ ▢ ▢ ▢ X", ma la quinta riga è "▢ ▢ X ▢ X". Puoi confermare se nella tua griglia la quinta riga è davvero uguale alla quarta, o se c'è una differenza nella posizione della "X" nell'ultima riga?Inizia la tua risposta con "RISPOSTA:" e non aggiungere altro.} \\
            }
        }
    }
    & \\ \\

    \theutterance \stepcounter{utterance}  
    & \multicolumn{4}{p{0.6\linewidth}}{
        \cellcolor[rgb]{0.9,0.9,0.9}{
            \makecell[{{p{\linewidth}}}]{
                \texttt{\tiny{[P1$\rangle$GM]}}
                \texttt{RISPOSTA: Nella mia griglia, la quinta riga è effettivamente identica alla quarta, composta da quattro "▢" seguiti da una "X" (▢ ▢ ▢ ▢ X). Non c'è una "X" nella terza posizione dell'ultima riga.} \\
            }
        }
    }
    & & \\ \\

    \theutterance \stepcounter{utterance}  
    & & & \multicolumn{2}{p{0.3\linewidth}}{
        \cellcolor[rgb]{0.9,0.9,0.9}{
            \makecell[{{p{\linewidth}}}]{
                \texttt{\tiny{[GM$|$GM]}}
                \texttt{continue} \\
            }
        }
    }
    & & \\ \\

    \theutterance \stepcounter{utterance}  
    & & \multicolumn{4}{p{0.6\linewidth}}{
        \cellcolor[rgb]{0.9,0.9,0.9}{
            \makecell[{{p{\linewidth}}}]{
                \texttt{\tiny{[P1$\langle$GM]}}
                \texttt{Ora fai una domanda per scoprire nuovi aspetti della mia griglia che potrebbero essere diversi dalla tua griglia. Inizia con "DOMANDA:" e non aggiungere altro.} \\
            }
        }
    }
    & \\ \\

    \theutterance \stepcounter{utterance}  
    & \multicolumn{4}{p{0.6\linewidth}}{
        \cellcolor[rgb]{0.9,0.9,0.9}{
            \makecell[{{p{\linewidth}}}]{
                \texttt{\tiny{[P1$\rangle$GM]}}
                \texttt{DOMANDA: Nella tua terza riga, hai indicato "▢ X X X X". Nella mia griglia, la terza riga è "▢ ▢ X X X". Puoi confermare se nella tua griglia la seconda colonna della terza riga è una "X" o un "▢"?} \\
            }
        }
    }
    & & \\ \\

    \theutterance \stepcounter{utterance}  
    & & & \multicolumn{2}{p{0.3\linewidth}}{
        \cellcolor[rgb]{0.9,0.9,0.9}{
            \makecell[{{p{\linewidth}}}]{
                \texttt{\tiny{[GM$|$GM]}}
                \texttt{continue} \\
            }
        }
    }
    & & \\ \\

    \theutterance \stepcounter{utterance}  
    & & \multicolumn{4}{p{0.6\linewidth}}{
        \cellcolor[rgb]{0.9,0.9,0.9}{
            \makecell[{{p{\linewidth}}}]{
                \texttt{\tiny{[GM$\rangle$P2]}}
                \texttt{RISPOSTA: Nella mia griglia, la quinta riga è effettivamente identica alla quarta, composta da quattro "▢" seguiti da una "X" (▢ ▢ ▢ ▢ X). Non c'è una "X" nella terza posizione dell'ultima riga.} \\
\texttt{DOMANDA: Nella tua terza riga, hai indicato "▢ X X X X". Nella mia griglia, la terza riga è "▢ ▢ X X X". Puoi confermare se nella tua griglia la seconda colonna della terza riga è una "X" o un "▢"?Inizia la tua risposta con "RISPOSTA:" e non aggiungere altro.} \\
            }
        }
    }
    & \\ \\

    \theutterance \stepcounter{utterance}  
    & & & \multicolumn{4}{p{0.6\linewidth}}{
        \cellcolor[rgb]{0.9,0.9,0.9}{
            \makecell[{{p{\linewidth}}}]{
                \texttt{\tiny{[GM$\langle$P2]}}
                \texttt{RISPOSTA: Nella mia griglia, la terza riga è "▢ X X X X". Quindi, la seconda colonna della terza riga è una "X".} \\
            }
        }
    }
     \\ \\

    \theutterance \stepcounter{utterance}  
    & & & \multicolumn{2}{p{0.3\linewidth}}{
        \cellcolor[rgb]{0.9,0.9,0.9}{
            \makecell[{{p{\linewidth}}}]{
                \texttt{\tiny{[GM$|$GM]}}
                \texttt{continue} \\
            }
        }
    }
    & & \\ \\

    \theutterance \stepcounter{utterance}  
    & & \multicolumn{4}{p{0.6\linewidth}}{
        \cellcolor[rgb]{0.9,0.9,0.9}{
            \makecell[{{p{\linewidth}}}]{
                \texttt{\tiny{[GM$\rangle$P2]}}
                \texttt{Ora fai una domanda per scoprire nuovi aspetti della mia griglia che potrebbero essere diversi dalla tua griglia. Inizia con "DOMANDA:" e non aggiungere altro.} \\
            }
        }
    }
    & \\ \\

    \theutterance \stepcounter{utterance}  
    & & & \multicolumn{4}{p{0.6\linewidth}}{
        \cellcolor[rgb]{0.9,0.9,0.9}{
            \makecell[{{p{\linewidth}}}]{
                \texttt{\tiny{[GM$\langle$P2]}}
                \texttt{DOMANDA: Nella mia griglia, la seconda riga è "▢ ▢ ▢ ▢ X". Nella tua descrizione precedente, hai confermato che la tua seconda riga ha quattro "▢" seguiti da una "X". Per essere assolutamente sicuri, puoi confermare che nella tua griglia la seconda riga è esattamente "▢ ▢ ▢ ▢ X" e non contiene nessuna "X" nelle prime quattro posizioni?} \\
            }
        }
    }
     \\ \\

    \theutterance \stepcounter{utterance}  
    & & & \multicolumn{2}{p{0.3\linewidth}}{
        \cellcolor[rgb]{0.9,0.9,0.9}{
            \makecell[{{p{\linewidth}}}]{
                \texttt{\tiny{[GM$|$GM]}}
                \texttt{continue} \\
            }
        }
    }
    & & \\ \\

    \theutterance \stepcounter{utterance}  
    & & \multicolumn{4}{p{0.6\linewidth}}{
        \cellcolor[rgb]{0.9,0.9,0.9}{
            \makecell[{{p{\linewidth}}}]{
                \texttt{\tiny{[P1$\langle$GM]}}
                \texttt{RISPOSTA: Nella mia griglia, la terza riga è "▢ X X X X". Quindi, la seconda colonna della terza riga è una "X".} \\
\texttt{DOMANDA: Nella mia griglia, la seconda riga è "▢ ▢ ▢ ▢ X". Nella tua descrizione precedente, hai confermato che la tua seconda riga ha quattro "▢" seguiti da una "X". Per essere assolutamente sicuri, puoi confermare che nella tua griglia la seconda riga è esattamente "▢ ▢ ▢ ▢ X" e non contiene nessuna "X" nelle prime quattro posizioni?Inizia la tua risposta con "RISPOSTA:" e non aggiungere altro.} \\
            }
        }
    }
    & \\ \\

    \theutterance \stepcounter{utterance}  
    & \multicolumn{4}{p{0.6\linewidth}}{
        \cellcolor[rgb]{0.9,0.9,0.9}{
            \makecell[{{p{\linewidth}}}]{
                \texttt{\tiny{[P1$\rangle$GM]}}
                \texttt{RISPOSTA: Sì, confermo che nella mia griglia la seconda riga è esattamente "▢ ▢ ▢ ▢ X" e non contiene nessuna "X" nelle prime quattro posizioni.} \\
            }
        }
    }
    & & \\ \\

    \theutterance \stepcounter{utterance}  
    & & & \multicolumn{2}{p{0.3\linewidth}}{
        \cellcolor[rgb]{0.9,0.9,0.9}{
            \makecell[{{p{\linewidth}}}]{
                \texttt{\tiny{[GM$|$GM]}}
                \texttt{continue} \\
            }
        }
    }
    & & \\ \\

    \theutterance \stepcounter{utterance}  
    & & \multicolumn{4}{p{0.6\linewidth}}{
        \cellcolor[rgb]{0.9,0.9,0.9}{
            \makecell[{{p{\linewidth}}}]{
                \texttt{\tiny{[P1$\langle$GM]}}
                \texttt{Ora fai una domanda per scoprire nuovi aspetti della mia griglia che potrebbero essere diversi dalla tua griglia. Inizia con "DOMANDA:" e non aggiungere altro.} \\
            }
        }
    }
    & \\ \\

    \theutterance \stepcounter{utterance}  
    & \multicolumn{4}{p{0.6\linewidth}}{
        \cellcolor[rgb]{0.9,0.9,0.9}{
            \makecell[{{p{\linewidth}}}]{
                \texttt{\tiny{[P1$\rangle$GM]}}
                \texttt{DOMANDA: Nella tua prima riga, hai indicato "X X X X X". Nella mia griglia, la prima riga è anch'essa "X X X X X". Per escludere qualsiasi differenza, puoi confermare che tutte e cinque le celle della prima riga sono effettivamente "X" e non ci sono "▢" in questa riga?} \\
            }
        }
    }
    & & \\ \\

    \theutterance \stepcounter{utterance}  
    & & & \multicolumn{2}{p{0.3\linewidth}}{
        \cellcolor[rgb]{0.9,0.9,0.9}{
            \makecell[{{p{\linewidth}}}]{
                \texttt{\tiny{[GM$|$GM]}}
                \texttt{continue} \\
            }
        }
    }
    & & \\ \\

    \theutterance \stepcounter{utterance}  
    & & \multicolumn{4}{p{0.6\linewidth}}{
        \cellcolor[rgb]{0.9,0.9,0.9}{
            \makecell[{{p{\linewidth}}}]{
                \texttt{\tiny{[GM$\rangle$P2]}}
                \texttt{RISPOSTA: Sì, confermo che nella mia griglia la seconda riga è esattamente "▢ ▢ ▢ ▢ X" e non contiene nessuna "X" nelle prime quattro posizioni.} \\
\texttt{DOMANDA: Nella tua prima riga, hai indicato "X X X X X". Nella mia griglia, la prima riga è anch'essa "X X X X X". Per escludere qualsiasi differenza, puoi confermare che tutte e cinque le celle della prima riga sono effettivamente "X" e non ci sono "▢" in questa riga?Inizia la tua risposta con "RISPOSTA:" e non aggiungere altro.} \\
            }
        }
    }
    & \\ \\

    \theutterance \stepcounter{utterance}  
    & & & \multicolumn{4}{p{0.6\linewidth}}{
        \cellcolor[rgb]{0.9,0.9,0.9}{
            \makecell[{{p{\linewidth}}}]{
                \texttt{\tiny{[GM$\langle$P2]}}
                \texttt{RISPOSTA: Sì, confermo che nella mia griglia la prima riga è esattamente "X X X X X" e tutte e cinque le celle sono "X".} \\
            }
        }
    }
     \\ \\

    \theutterance \stepcounter{utterance}  
    & & & \multicolumn{2}{p{0.3\linewidth}}{
        \cellcolor[rgb]{0.9,0.9,0.9}{
            \makecell[{{p{\linewidth}}}]{
                \texttt{\tiny{[GM$|$GM]}}
                \texttt{continue} \\
            }
        }
    }
    & & \\ \\

    \theutterance \stepcounter{utterance}  
    & & \multicolumn{4}{p{0.6\linewidth}}{
        \cellcolor[rgb]{0.9,0.9,0.9}{
            \makecell[{{p{\linewidth}}}]{
                \texttt{\tiny{[GM$\rangle$P2]}}
                \texttt{Ora fai una domanda per scoprire nuovi aspetti della mia griglia che potrebbero essere diversi dalla tua griglia. Inizia con "DOMANDA:" e non aggiungere altro.} \\
            }
        }
    }
    & \\ \\

    \theutterance \stepcounter{utterance}  
    & & & \multicolumn{4}{p{0.6\linewidth}}{
        \cellcolor[rgb]{0.9,0.9,0.9}{
            \makecell[{{p{\linewidth}}}]{
                \texttt{\tiny{[GM$\langle$P2]}}
                \texttt{DOMANDA: Nella mia griglia, la terza riga è "▢ X X X X". Nella tua descrizione iniziale, hai indicato che la terza riga è "▢ X X X X". Per essere sicuri che non ci siano differenze nelle colonne successive, puoi confermare che nella tua griglia la terza riga ha esattamente tre "X" consecutive dopo il primo "▢", ovvero le celle nella colonna 2, 3, 4 e 5 della terza riga sono tutte "X"?} \\
            }
        }
    }
     \\ \\

    \theutterance \stepcounter{utterance}  
    & & & \multicolumn{2}{p{0.3\linewidth}}{
        \cellcolor[rgb]{0.9,0.9,0.9}{
            \makecell[{{p{\linewidth}}}]{
                \texttt{\tiny{[GM$|$GM]}}
                \texttt{continue} \\
            }
        }
    }
    & & \\ \\

    \theutterance \stepcounter{utterance}  
    & & \multicolumn{4}{p{0.6\linewidth}}{
        \cellcolor[rgb]{0.9,0.9,0.9}{
            \makecell[{{p{\linewidth}}}]{
                \texttt{\tiny{[P1$\langle$GM]}}
                \texttt{RISPOSTA: Sì, confermo che nella mia griglia la prima riga è esattamente "X X X X X" e tutte e cinque le celle sono "X".} \\
\texttt{DOMANDA: Nella mia griglia, la terza riga è "▢ X X X X". Nella tua descrizione iniziale, hai indicato che la terza riga è "▢ X X X X". Per essere sicuri che non ci siano differenze nelle colonne successive, puoi confermare che nella tua griglia la terza riga ha esattamente tre "X" consecutive dopo il primo "▢", ovvero le celle nella colonna 2, 3, 4 e 5 della terza riga sono tutte "X"?Inizia la tua risposta con "RISPOSTA:" e non aggiungere altro.} \\
            }
        }
    }
    & \\ \\

    \theutterance \stepcounter{utterance}  
    & \multicolumn{4}{p{0.6\linewidth}}{
        \cellcolor[rgb]{0.9,0.9,0.9}{
            \makecell[{{p{\linewidth}}}]{
                \texttt{\tiny{[P1$\rangle$GM]}}
                \texttt{RISPOSTA: No, nella mia griglia la terza riga è "▢ ▢ X X X". Quindi, dopo il primo "▢", c'è un altro "▢" nella seconda colonna, seguito da tre "X" nelle colonne 3, 4 e 5. Le celle nella colonna 2 della terza riga sono "▢", non "X".} \\
            }
        }
    }
    & & \\ \\

    \theutterance \stepcounter{utterance}  
    & & & \multicolumn{2}{p{0.3\linewidth}}{
        \cellcolor[rgb]{0.9,0.9,0.9}{
            \makecell[{{p{\linewidth}}}]{
                \texttt{\tiny{[GM$|$GM]}}
                \texttt{continue} \\
            }
        }
    }
    & & \\ \\

    \theutterance \stepcounter{utterance}  
    & & \multicolumn{4}{p{0.6\linewidth}}{
        \cellcolor[rgb]{0.9,0.9,0.9}{
            \makecell[{{p{\linewidth}}}]{
                \texttt{\tiny{[P1$\langle$GM]}}
                \texttt{Ora fai una domanda per scoprire nuovi aspetti della mia griglia che potrebbero essere diversi dalla tua griglia. Inizia con "DOMANDA:" e non aggiungere altro.} \\
            }
        }
    }
    & \\ \\

    \theutterance \stepcounter{utterance}  
    & \multicolumn{4}{p{0.6\linewidth}}{
        \cellcolor[rgb]{0.9,0.9,0.9}{
            \makecell[{{p{\linewidth}}}]{
                \texttt{\tiny{[P1$\rangle$GM]}}
                \texttt{DOMANDA: Nella tua griglia, la terza riga è "▢ X X X X". Nella mia griglia, la terza riga è "▢ ▢ X X X". Puoi confermare se nella tua griglia la seconda colonna della terza riga è effettivamente una "X" o un "▢"?} \\
            }
        }
    }
    & & \\ \\

    \theutterance \stepcounter{utterance}  
    & & & \multicolumn{2}{p{0.3\linewidth}}{
        \cellcolor[rgb]{0.9,0.9,0.9}{
            \makecell[{{p{\linewidth}}}]{
                \texttt{\tiny{[GM$|$GM]}}
                \texttt{continue} \\
            }
        }
    }
    & & \\ \\

    \theutterance \stepcounter{utterance}  
    & & \multicolumn{4}{p{0.6\linewidth}}{
        \cellcolor[rgb]{0.9,0.9,0.9}{
            \makecell[{{p{\linewidth}}}]{
                \texttt{\tiny{[GM$\rangle$P2]}}
                \texttt{RISPOSTA: No, nella mia griglia la terza riga è "▢ ▢ X X X". Quindi, dopo il primo "▢", c'è un altro "▢" nella seconda colonna, seguito da tre "X" nelle colonne 3, 4 e 5. Le celle nella colonna 2 della terza riga sono "▢", non "X".} \\
\texttt{DOMANDA: Nella tua griglia, la terza riga è "▢ X X X X". Nella mia griglia, la terza riga è "▢ ▢ X X X". Puoi confermare se nella tua griglia la seconda colonna della terza riga è effettivamente una "X" o un "▢"?Inizia la tua risposta con "RISPOSTA:" e non aggiungere altro.} \\
            }
        }
    }
    & \\ \\

    \theutterance \stepcounter{utterance}  
    & & & \multicolumn{4}{p{0.6\linewidth}}{
        \cellcolor[rgb]{0.9,0.9,0.9}{
            \makecell[{{p{\linewidth}}}]{
                \texttt{\tiny{[GM$\langle$P2]}}
                \texttt{RISPOSTA: Nella mia griglia, la terza riga è "▢ X X X X". Quindi, la seconda colonna della terza riga è una "X".} \\
            }
        }
    }
     \\ \\

    \theutterance \stepcounter{utterance}  
    & & & \multicolumn{2}{p{0.3\linewidth}}{
        \cellcolor[rgb]{0.9,0.9,0.9}{
            \makecell[{{p{\linewidth}}}]{
                \texttt{\tiny{[GM$|$GM]}}
                \texttt{continue} \\
            }
        }
    }
    & & \\ \\

    \theutterance \stepcounter{utterance}  
    & & \multicolumn{4}{p{0.6\linewidth}}{
        \cellcolor[rgb]{0.9,0.9,0.9}{
            \makecell[{{p{\linewidth}}}]{
                \texttt{\tiny{[GM$\rangle$P2]}}
                \texttt{Ora prendi una decisione. Cosa ne pensi: la tua griglia e l'altra griglia descritta sono uguali?  Scrivi "DECISIONE: stessa griglia" se pensi che siano la stessa griglia oppure "DECISIONE: griglie diverse" se pensi che siano griglie diverse. Non aggiungere nient'altro.} \\
            }
        }
    }
    & \\ \\

    \theutterance \stepcounter{utterance}  
    & & & \multicolumn{4}{p{0.6\linewidth}}{
        \cellcolor[rgb]{0.9,0.9,0.9}{
            \makecell[{{p{\linewidth}}}]{
                \texttt{\tiny{[GM$\langle$P2]}}
                \texttt{DECISIONE: griglie diverse} \\
            }
        }
    }
     \\ \\

    \theutterance \stepcounter{utterance}  
    & & & \multicolumn{2}{p{0.3\linewidth}}{
        \cellcolor[rgb]{0.9,0.9,0.9}{
            \makecell[{{p{\linewidth}}}]{
                \texttt{\tiny{[GM$|$GM]}}
                \texttt{continue} \\
            }
        }
    }
    & & \\ \\

    \theutterance \stepcounter{utterance}  
    & & & \multicolumn{2}{p{0.3\linewidth}}{
        \cellcolor[rgb]{0.9,0.9,0.9}{
            \makecell[{{p{\linewidth}}}]{
                \texttt{\tiny{[GM$|$GM]}}
                \texttt{success} \\
            }
        }
    }
    & & \\ \\

    \theutterance \stepcounter{utterance}  
    & & \multicolumn{4}{p{0.6\linewidth}}{
        \cellcolor[rgb]{0.9,0.9,0.9}{
            \makecell[{{p{\linewidth}}}]{
                \texttt{\tiny{[P1$\langle$GM]}}
                \texttt{RISPOSTA: Nella mia griglia, la terza riga è "▢ X X X X". Quindi, la seconda colonna della terza riga è una "X".} \\
\texttt{Ora prendi una decisione. Cosa ne pensi: la tua griglia e l'altra griglia descritta sono uguali?  Scrivi "DECISIONE: stessa griglia" se pensi che siano la stessa griglia oppure "DECISIONE: griglie diverse" se pensi che siano griglie diverse. Non aggiungere nient'altro.} \\
            }
        }
    }
    & \\ \\

    \theutterance \stepcounter{utterance}  
    & \multicolumn{4}{p{0.6\linewidth}}{
        \cellcolor[rgb]{0.9,0.9,0.9}{
            \makecell[{{p{\linewidth}}}]{
                \texttt{\tiny{[P1$\rangle$GM]}}
                \texttt{DECISIONE: griglie diverse} \\
            }
        }
    }
    & & \\ \\

    \theutterance \stepcounter{utterance}  
    & & & \multicolumn{2}{p{0.3\linewidth}}{
        \cellcolor[rgb]{0.9,0.9,0.9}{
            \makecell[{{p{\linewidth}}}]{
                \texttt{\tiny{[GM$|$GM]}}
                \texttt{continue} \\
            }
        }
    }
    & & \\ \\

    \theutterance \stepcounter{utterance}  
    & & & \multicolumn{2}{p{0.3\linewidth}}{
        \cellcolor[rgb]{0.9,0.9,0.9}{
            \makecell[{{p{\linewidth}}}]{
                \texttt{\tiny{[GM$|$GM]}}
                \texttt{success} \\
            }
        }
    }
    & & \\ \\

\end{supertabular}
}

\end{document}
